\section{Discussion, recommendations and limitations}\label{sec:discussion}

\paragraph{What the evidence supports.} On the collections we examined, rules that test a model of the conditional distribution do not beat rules that read marginal asymmetry or error scale, and no rule leads across collections. This is not because the dependence-based rules are poor: ANM-HSIC, RESIT and the Gaussian-process ANM are right on $91$ to $97\%$ of the least heteroscedastic third of the T\"ubingen pairs and ANM-HSIC is perfect in simulation when its assumption holds. They are reversed when the noise scale depends on the cause, degraded by measurement noise, and unable to use the marginal cue that orients the chamber collections; averaged over a collection that mixes regimes they equal the baselines. The baselines are not robust either: IGCI is at $8\%$ on the Malus pairs and, with or without a jitter on ties, at $6\%$ or $94\%$ on the i.i.d.\ chambers; RECI is at $25\%$ on the Bernoulli pairs; compression is at $38\%$ on the i.i.d.\ chambers. We do not claim that published methods misreport (we reproduced the one table we audited), nor that dependence-based rules cannot be improved: the failed regime-gated rule and the failed predictive validation show the account is incomplete.

\paragraph{Recommendations.}
\begin{enumerate}
\item \emph{Test paired and state the power.} With $102$ pairs the detectable difference is about $19$ points ($27$ after correcting for $78$ comparisons); a gain of two to five points is not resolvable, and about $330$ pairs are needed for $10$ points.
\item \emph{Report the coverage curve and both weightings.} A single decided-only, weighted or unweighted number selects one point of a multiverse that moves a rule by $11$ to $26$ points and its rank by up to nine places.
\item \emph{Fix preprocessing and settings before the data are seen and publish the code,} including conventions the text does not fix: a rounding rule moved a weighted headline by six points, IGCI's reference measure moved its accuracy by $14$ points, and the treatment of ties turned $94\%$ into $6\%$.
\item \emph{Report several collections and test a marginal-only classifier,} which is cheap and shows how much of a benchmark can be solved without dependence structure; rankings on one collection say little about another.
\item \emph{Report the effective sample size and the number of distinct edges} (seventeen of the nineteen untouched chamber recordings have an effective sample size of about three out of $4000$ rows; the new pair collections have four to eight distinct edges).
\item \emph{For learned methods, show real-data accuracy next to the simulated test} and how far the real diagnostics lie from the simulator's support.
\end{enumerate}

\subsection{Limitations}\label{sec:limits}
T\"ubingen has $102$ pairs, Sachs $17$ arcs and the other collections four to $47$ distinct edges, so all inference is as imprecise as Section~\ref{sec:results} says. IGCI, RECI, ANM-HSIC, pairwise LiNGAM, compression and TRA are our implementations (RECI agrees with a published re-run within $0.2$ weighted points; TRA uses our constants); we did not run SLOPE, CAM, SCORE or PNL, and tuning could change a ranking. The multiverse is finite, its choices are ours, its accuracies are computed on overlapping pairs and several scores are not confidences. The heteroscedasticity measure is computed in the true direction and the tercile analysis is post hoc; the manipulation and factorial map use one family of mechanisms, and the factorial simulator does not cover the real collections (Appendix~\ref{app:mech}). The marginal-only classifier is small, and its chamber results concern an actuator-to-sensor design. Why additive-noise rules lose on the chambers and why IGCI fails on Malus are not diagnosed. Analyses designed after the data are labelled as such, and several errors found on internal review (leakage in the marginal classifier, a confounded manipulation, universes with too few pairs) are described where they occurred.
